\section{Représentation des programmes}
\subsection{Arbres de syntaxe abstraits}

\begin{frame}[fragile,t]
    \frametitle{Arbres de syntaxe abstraits (AST)}

    \begin{minted}{ocaml}
let rec sum (l : int list) : int =
    match l with
    | [] -> 0
    | x :: q -> x + sum q
    \end{minted}

    \maxsizebox{\textwidth}{0.6\textheight}{
        \begin{forest}
            [{\texttt{FunctionDefinition(isRec: true, name: sum)}}
            [{body: \texttt{Match}},s sep=3cm
            [value : \texttt{Variable(l)}]
            [cases,s sep=3cm
            [\texttt{MatchCase}
            [{pattern : \texttt{EmptyList}},tier=forAlignment]
            [{expression : \texttt{Integer(0)}},tier=forAlignment]
            ]
            [\texttt{MatchCase},tier=forAlignment
            [pattern : \texttt{Concatenation}[...]]
            [expression : \texttt{Addition}
            [lhs : \texttt{Variable(x)}]
            [rhs : ...]
            ]
            ]
            ]
            ]
            [parameters: ...]
            [returnType: ...]
            ]
        \end{forest}
    }
\end{frame}

\begin{frame}[fragile,t]
    \frametitle{AST : Limitations}

    Les AST sont peu adaptés aux algorithmes demandant une connaissance de
    l'ordre d'évaluation du programme.

    \vspace{0.5cm}

    \begin{minipage}[t]{0.5\textwidth}
        \begin{minted}{ocaml}
a + b
        \end{minted}

        \begin{forest}
            [\texttt{Addition}
                [\texttt{Variable(a)}]
                [\texttt{Variable(b)}]]
        \end{forest}

        \only<2->{
            \begin{enumerate}
                \item \texttt{b}
                \item \texttt{a}
                \item \texttt{Addition}
            \end{enumerate}
        }
    \end{minipage}%
    \begin{minipage}[t]{0.5\textwidth}
        \begin{minted}{ocaml}
if a then b
        \end{minted}

        \begin{forest}
            [\texttt{If}
                [\texttt{Variable(a)}]
                [\texttt{Variable(b)}]]
        \end{forest}

        \only<2->{
            \begin{enumerate}
                \item \texttt{a}
                \item (Peut-être) \texttt{b}
            \end{enumerate}
        }
    \end{minipage}

    \vspace{0.5cm}

    \only<3->{
        \textcolor{orange}{$\blacktriangleright$} On préférerait une structure ne nécessitant pas de discriminer selon les étiquettes des nœuds.
    }
\end{frame}

\subsection{Formes linéaires}

\begin{frame}[fragile]
    \frametitle{Expressions localement terminales}

    \begin{definition}
        Une sous-expression $e'$ d'une expression \texttt{OCaml} $e$ est dite
        \textit{localement terminale} si l'évaluation de $e$ se termine
        immédiatement après l'évaluation de $e'$, à moins que $e'$ ne soit pas
        évaluée pour évaluer $e$.
    \end{definition}

    \pause

    \begin{minipage}[t]{0.5\textwidth}
        \begin{minted}{ocaml}
if est_pair n then
    n / 2
else
    (n - 1) / 2
        \end{minted}
    \end{minipage}%
    \begin{minipage}[t]{0.5\textwidth}
        \begin{minted}{ocaml}
let x = fibonacci 50 in
x * 10
        \end{minted}
    \end{minipage}

\end{frame}

\begin{frame}[t]
    \frametitle{Opérations élémentaires}
    \small
    \vspace{-0.2cm}

    \begin{definition}
        Si $C(e_1, \ldots, e_n)$ est une règle de construction d'une expression
        \texttt{OCaml}, alors on peut créer une \textit{opération élémentaire}
        avec la règle $C(a_1, \ldots, a_k, l_1, \ldots, l_{n - k})$, o\`u:

        \begin{itemize}
            \item Les $a_i$ sont des noms de variables pour chaque
                  sous-expression non localement terminale;
            \item Les $l_i$ sont des formes linéaires pour chaque
                  sous-expression localement terminale.
        \end{itemize}
    \end{definition}
\end{frame}

\addtocounter{framenumber}{-1}
\begin{frame}[t]
    \frametitle{Opérations élémentaires}
    \small
    \vspace{-0.2cm}

    \begin{definition}
        Si $C(e_1, \ldots, e_n)$ est une règle de construction d'une expression
        \texttt{OCaml}, alors on peut créer une \textit{opération élémentaire}
        avec la règle $C(a_1, \ldots, a_k, l_1, \ldots, l_{n - k})$.
    \end{definition}

    $$
        \texttt{If}_{\texttt{OCaml}}: \left\{ \begin{aligned}
            E \times E \times E & \rightarrow E                                                            \\
            (e_1, e_2, e_3)     & \rightarrowtail \texttt{if } e_1 \texttt{ then } e_2 \texttt{ else } e_3
        \end{aligned} \right.
    $$
    \pause
    $$
        \texttt{If}_{\text{Opération élémentaire}}: \left\{ \begin{aligned}
            V \times L \times L & \rightarrow O                                                            \\
            (a_1, l_1, l_2)     & \rightarrowtail \texttt{if } a_1 \texttt{ then } l_1 \texttt{ else } l_2
        \end{aligned} \right.
    $$

    Avec:
    \begin{itemize}
        \item $E$ l'ensemble des expressions \texttt{OCaml};
        \item $V$ l'ensemble des noms de variable \texttt{OCaml};
        \item $L$ l'ensemble des formes linéaires;
        \item $O$ l'ensemble des opérations élémentaires.
    \end{itemize}
\end{frame}


\begin{frame}
    \frametitle{Opérations élémentaires : exemples}

    \begin{minipage}{0.5\textwidth}
        $$
            \texttt{+} : \left\{
            \begin{aligned}
                V \times V & \rightarrow O                        \\
                (a_1, a_2) & \rightarrowtail a_1 \texttt{ + } a_2
            \end{aligned} \right.
        $$
    \end{minipage}%
    \begin{minipage}{0.5\textwidth}
        \texttt{a + b} \\
        \texttt{$a_1$ + mon\_super\_identifiant}
    \end{minipage}

    \pause

    \begin{minipage}{0.6\textwidth}
        $$
            \text{Tuple} : \left\{
            \begin{aligned}
                V^k                & \rightarrow O                                                                \\
                (a_1, \ldots, a_k) & \rightarrowtail \texttt{(} a_1 \texttt{, } \ldots \texttt{, } a_l \texttt{)}
            \end{aligned} \right.
        $$
    \end{minipage}%
    \begin{minipage}{0.4\textwidth}
        \texttt{(a, b, c)} \\
        \texttt{(a, $a_2$, $a_3$, d)}
    \end{minipage}

\end{frame}

\begin{frame}[t]
    \frametitle{Formes linéaires}

    \begin{definition}
        Une \textit{forme linéaire} est une liste non vide de couples $(a, e)$ avec $a$
        un nom de variable \texttt{OCaml} et $e$ une opération élémentaire.
    \end{definition}

    \pause

    On notera la forme linéaire $[(a_1, e_1), (a_2, e_2), \ldots]$ comme suit :
    \begin{align*}
        a_1 & := e_1 \\
        a_2 & := e_2 \\
        \ldots
    \end{align*}
\end{frame}

\begin{frame}
    \frametitle{Formes linéaires : exemples}

    \begin{align*}
        a_1 & := \texttt{a + b}                \\
        a_2 & := \texttt{$a_1$ + c}            \\
        a_3 & := \texttt{(a, b, c)}            \\
        a_4 & := \texttt{(a, $a_2$, $a_3$, d)}
    \end{align*}

    \vspace{0.5cm}

    \pause

    Évaluation avec $\texttt{a} = 3$, $\texttt{b} = 2$, $\texttt{c} = 10$ et
    $\texttt{d} = -1$ :
    \begin{align*}
        a_1 & = 5                       \\
        a_2 & = 15                      \\
        a_3 & = (3, 2, 10)              \\
        a_4 & = (3, 15, (3, 2, 10), -1)
    \end{align*}
\end{frame}

\addtocounter{framenumber}{-1}
\begin{frame}[fragile]
    \frametitle{Formes linéaires : exemples}

    \begin{align*}
        a_1 := & \ \texttt{true}                              \\
        a_2 := & \ \texttt{if } a_1 \texttt{ then}            \\
               & \hspace{0.5cm} \begin{aligned}
                                    a_3 & := \texttt{6}           \\
                                    a_4 & := \texttt{7}           \\
                                    a_5 & := a_3 \texttt{ * } a_4
                                \end{aligned} \\
               & \ \texttt{else}                              \\
               & \hspace{0.5cm} \begin{aligned}
                                    a_6 & := \texttt{17}
                                \end{aligned}
    \end{align*}

    \pause

    \begin{minted}{ocaml}
if true then
    6 * 7
else
    17
    \end{minted}
\end{frame}

\subsection{Linéarisation et délinéarisation}

\begin{frame}[t]
    \frametitle{Linéarisation}

    On définit inductivement la linéarisation :

    \begin{itemize}
        \pause
        \item Linéariser(\texttt{a + b}) = Linéariser(\texttt{b}).Linéariser(\texttt{a}).[$a_k$ := $a_i$ \texttt{+} $a_j$] \\
              \pause
        \item Linéariser(\texttt{if a then b}) = \\
              \hspace{0.5cm} Linéariser(\texttt{a}).[$a_k$ := \texttt{if} $a_i$ \texttt{then} Linéariser(\texttt{b})]
              \pause
        \item ...
    \end{itemize}
    \small \textit{Avec $a_k$, $a_i$ et $a_j$ choisis en fonction des noms déjà utilisés}

    \pause

    \vspace{1cm}

    \normalsize

    \begin{itemize}
        \item On démontre que la linéarisation préserve le comportement du
              programme.
    \end{itemize}
\end{frame}

\begin{frame}[fragile,t]
    \frametitle{Délinéarisation}

    On peut aussi définir un algorithme transformant une forme linéaire en un
    programme \texttt{OCaml}:

    $$a_k := expr \longrightarrow \texttt{let a\_k = expr}$$

    \vspace{0.5cm}
    \pause

    \begin{minipage}{0.5\textwidth}
        Délinéariser\small$\left(\begin{aligned}
                    a_1 & := \texttt{a}                                \\
                    a_2 & := \texttt{2}                                \\
                    a_3 & := a_1 \texttt{ = } a_2                      \\
                    a_4 & := \texttt{if } a_3 \texttt{ then}           \\
                        & \hspace{0.5cm} \begin{aligned}
                                             a_5 & := 3                    \\
                                             a_6 & := b                    \\
                                             a_7 & := a_6 \texttt{ * } a_5
                                         \end{aligned} \\
                        & \texttt{else } ...
                \end{aligned}\right)$
    \end{minipage}=%
    \begin{minipage}{0.45\textwidth}
        \small \begin{minted}{ocaml}
let a_1 = a in
let a_2 = 2 in
let a_3 = a_1 = a_2 in
if a_3 then
    let a_5 = 3 in
    let a_6 = b in
    a_6 * a_5
else
    ...
            \end{minted}
    \end{minipage}
\end{frame}

\addtocounter{framenumber}{-1}
\begin{frame}[fragile,t]
    \frametitle{Délinéarisation}

    On peut aussi définir un algorithme transformant une forme linéaire en un
    programme \texttt{OCaml}:

    $$a_k := expr \longrightarrow \texttt{let a\_k = expr}$$

    \vspace{0.5cm}

    \begin{minipage}{0.5\textwidth}
        Délinéariser\small$\left(\begin{aligned}
                    a_1 & := \texttt{a}                                \\
                    a_2 & := \texttt{2}                                \\
                    a_3 & := a_1 \texttt{ = } a_2                      \\
                    a_4 & := \texttt{if } a_3 \texttt{ then}           \\
                        & \hspace{0.5cm} \begin{aligned}
                                             a_5 & := 3                    \\
                                             a_6 & := b                    \\
                                             a_7 & := a_6 \texttt{ * } a_5
                                         \end{aligned} \\
                        & \texttt{else } ...
                \end{aligned}\right)$
    \end{minipage}=%
    \begin{minipage}{0.45\textwidth}
        \small \begin{minted}{ocaml}
if a = 2 then
    b * 3
else
    a - 2
            \end{minted}
    \end{minipage}
\end{frame}

